\begin{summarybox}
\textbf{\textcolor{CSummary}{一句话核心}}
\textbf{诱导公式通过单位圆对称性把任意角三角函数化为参考角三角函数；求特殊值时再进一步化为锐角三角函数。口诀“奇变偶不变，符号看象限”用于快速判断函数名与符号。}

\medskip
\textbf{\textcolor{CSummary}{使用条件}}
\begin{itemize}[leftmargin=*, itemsep=0.5em, topsep=0.4em]
    \item \textbf{条件 1（角度形式）：} 常将角写成 $\frac{k\pi}{2}\pm\alpha$，其中 $k\in\mathbb{Z}$。
    \item \textbf{条件 2（任意角与参考角）：} 公式中的 $\alpha$ 可以是任意角；口诀判号时，才把 $\alpha$ 暂看作锐角或参考角。
    \item \textbf{条件 3（定义域）：} 涉及 $\tan$、$\cot$ 时必须保证左右两边都有意义。$\tan\theta$ 要求 $\cos\theta\neq0$，$\cot\theta$ 要求 $\sin\theta\neq0$。
\end{itemize}

\medskip
\textbf{\textcolor{CSummary}{关键公式}}
\[
\begin{aligned}
\sin(\pi-\alpha)&=\sin\alpha, & \cos(\pi-\alpha)&=-\cos\alpha,\\
\sin(\pi+\alpha)&=-\sin\alpha, & \cos(\pi+\alpha)&=-\cos\alpha,\\
\tan(\pi-\alpha)&=-\tan\alpha, & \tan(\pi+\alpha)&=\tan\alpha.
\end{aligned}
\]
\[
\begin{aligned}
\sin\left(\frac{\pi}{2}-\alpha\right)&=\cos\alpha, &
\cos\left(\frac{\pi}{2}-\alpha\right)&=\sin\alpha,\\
\sin\left(\frac{\pi}{2}+\alpha\right)&=\cos\alpha, &
\cos\left(\frac{\pi}{2}+\alpha\right)&=-\sin\alpha.
\end{aligned}
\]
含 $\tan$、$\cot$ 的变名公式同样成立，但必须先检查定义域。

\medskip
\textbf{\textcolor{CSummary}{关键提醒}}
\begin{itemize}[leftmargin=*, itemsep=0.5em, topsep=0.4em]
    \item \textbf{符号判定：} 口诀判号时，将 $\alpha$ 暂看作参考角；不要把这句话误解为公式只对锐角成立。
    \item \textbf{k的奇偶：} 准确数出 $k$ 数值，如 $\pi=2\cdot\frac{\pi}{2}$，所以 $k=2$ 为偶数，函数名不变。
    \item \textbf{周期与偶性：} 例如 $\cos(2\pi-\alpha)=\cos(-\alpha)=\cos\alpha$，这里先用周期性，再用余弦的偶性。
\end{itemize}
\end{summarybox}
